\documentclass[preprint,12pt]{elsarticle}

\usepackage{amsmath,amssymb}
\usepackage{graphicx}
\usepackage{booktabs}
\usepackage{longtable}
\usepackage{multirow}
\usepackage{hyperref}
\usepackage{xcolor}
\usepackage[margin=2.5cm]{geometry}

\journal{Computers, Environment and Urban Systems}

\begin{document}

\begin{frontmatter}

\title{Supplementary Material for:\\
An Open-Source Pipeline for Spatiotemporal Traffic Congestion Analysis: Integrating HERE Traffic, OSMnx, and PySAL across Three Indonesian Cities}

\author[ugm]{Firman Hadi\corref{cor1}}
\ead{firman.hadi@mail.ugm.ac.id}
\author[undip]{Yasser Wahyuddin}
\author[undip]{L.M. Sabri}
\author[its]{Agung Indrajit}

\affiliation[ugm]{organization={Department of Geodetic Engineering, Faculty of Engineering, Universitas Gadjah Mada},
  city={Yogyakarta}, country={Indonesia}}
\affiliation[undip]{organization={Department of Geodetic Engineering, Faculty of Engineering, Universitas Diponegoro},
  city={Semarang}, country={Indonesia}}
\affiliation[its]{organization={Department of Geomatics Engineering, Faculty of Civil, Planning, and Geo Engineering, Institut Teknologi Sepuluh Nopember},
  city={Surabaya}, country={Indonesia}}

\cortext[cor1]{Corresponding author.}

\end{frontmatter}

\setcounter{section}{0}
\renewcommand{\thesection}{S\arabic{section}}
\renewcommand{\thetable}{S\arabic{table}}
\renewcommand{\thefigure}{S\arabic{figure}}

%% ============================================================
\section{Spatial Weight Sensitivity Analysis}
\label{sec:s1}

To confirm that the null global spatial autocorrelation finding reported in the main text is robust, we computed Moran's I under five spatial weight specifications for each city: K-nearest neighbours with $k \in \{4, 8, 12\}$ and distance-band weights at 500\,m and 1000\,m thresholds.

\begin{table}[!htbp]
\centering
\caption{Moran's I under alternative spatial weight specifications (overall mean jam factor). All 15 city--weight combinations yield non-significant results ($p > 0.05$), confirming the absence of global spatial autocorrelation.}
\label{tab:s1}
\begin{tabular}{lccccc}
\toprule
City & KNN $k{=}4$ & KNN $k{=}8$ & KNN $k{=}12$ & Dist 500\,m & Dist 1000\,m \\
\midrule
Jakarta  & 0.001 ($p{=}0.86$) & 0.003 ($p{=}0.46$) & 0.005 ($p{=}0.12$) & $-0.009$ ($p{=}0.14$) & $-0.005$ ($p{=}0.33$) \\
Bandung  & 0.004 ($p{=}0.74$) & 0.007 ($p{=}0.35$) & 0.009 ($p{=}0.20$) & 0.007 ($p{=}0.63$)    & $-0.000$ ($p{=}0.99$) \\
Semarang & 0.008 ($p{=}0.67$) & $-0.004$ ($p{=}0.84$) & 0.003 ($p{=}0.77$) & 0.003 ($p{=}0.89$) & $-0.004$ ($p{=}0.88$) \\
\bottomrule
\end{tabular}
\end{table}

The distance-band weights at 500\,m and 1000\,m provide implicit correlogram information. Both distance thresholds yield non-significant Moran's I across all cities, indicating that spatial autocorrelation is absent at both local (neighbouring segments) and broader spatial scales.

%% ============================================================
\section{Period-Specific Moran's I}
\label{sec:s2}

Global Moran's I was computed separately for each of the eight temporal periods to test whether spatial clustering emerges during specific times of day. Spatial weights use KNN with $k{=}8$.

\begin{table}[!htbp]
\centering
\caption{Moran's I by temporal period (KNN $k{=}8$). No period--city combination reaches conventional significance ($p < 0.05$), confirming that the absence of global spatial clustering is not an artefact of temporal averaging.}
\label{tab:s2}
\begin{tabular}{lccc}
\toprule
Period & Jakarta & Bandung & Semarang \\
\midrule
Night             & 0.000 ($p{=}0.91$)    & 0.002 ($p{=}0.81$)    & 0.004 ($p{=}0.74$)    \\
Morning Peak      & $-0.001$ ($p{=}0.85$) & 0.009 ($p{=}0.27$)    & $-0.003$ ($p{=}0.86$) \\
Morning Off-Peak  & 0.005 ($p{=}0.19$)    & 0.009 ($p{=}0.28$)    & $-0.008$ ($p{=}0.62$) \\
Lunch Hours       & 0.001 ($p{=}0.79$)    & 0.003 ($p{=}0.68$)    & $-0.022$ ($p{=}0.14$) \\
Afternoon Off-Peak& 0.006 ($p{=}0.10$)    & 0.016 ($p{=}0.058$)   & $-0.001$ ($p{=}0.99$) \\
Evening Peak      & 0.003 ($p{=}0.46$)    & 0.007 ($p{=}0.35$)    & $-0.004$ ($p{=}0.84$) \\
Evening Off-Peak  & 0.004 ($p{=}0.25$)    & $-0.002$ ($p{=}0.81$) & 0.001 ($p{=}0.89$)    \\
Late Night        & $-0.002$ ($p{=}0.61$) & 0.005 ($p{=}0.51$)    & $-0.018$ ($p{=}0.23$) \\
\bottomrule
\end{tabular}
\end{table}

The closest result to significance is Bandung's afternoon off-peak period ($I{=}0.016$, $p{=}0.058$), which does not survive conventional thresholds and would not survive Bonferroni correction for 24 simultaneous tests ($\alpha_{\text{adj}} = 0.05/24 = 0.002$).

%% ============================================================
\section{Getis-Ord Gi* Hot and Cold Spot Counts}
\label{sec:s3}

Getis-Ord Gi* statistics were computed with KNN ($k{=}8$) binary weights and 999 permutations to triangulate the LISA findings reported in the main text.

\begin{table}[!htbp]
\centering
\caption{Getis-Ord Gi* hot and cold spot counts (evening peak, $p < 0.05$). Cold spots outnumber hot spots in all three cities, indicating that clusters of low congestion are slightly more prevalent than clusters of high congestion.}
\label{tab:s3}
\begin{tabular}{lcccc}
\toprule
City & Hot Spots & Cold Spots & Not Significant & \% Significant \\
\midrule
Jakarta  & 660 (4.5\%) & 807 (5.5\%) & 13,082 (89.9\%) & 10.1\% \\
Bandung  & 118 (3.8\%) & 178 (5.8\%) & 2,773 (90.4\%)  & 9.6\%  \\
Semarang & 39 (3.6\%)  & 53 (4.9\%)  & 984 (91.4\%)    & 8.6\%  \\
\bottomrule
\end{tabular}
\end{table}

Among hot spots, 95 (14.4\%) in Jakarta, 15 (12.7\%) in Bandung, and 2 (5.1\%) in Semarang reach the stricter $p < 0.01$ threshold. The overall proportion of significant segments ($\sim$9--10\%) is consistent with the LISA results reported in the main text.

%% ============================================================
\section{Spatial Match Quality}
\label{sec:s4}

The spatial join between HERE Traffic segments and OSMnx road network edges was validated by examining centroid-to-centroid match distances.

\begin{table}[!htbp]
\centering
\caption{Spatial match quality: HERE Traffic segments to OSMnx edges. Median distances are under 20\,m in all cities, and the 200\,m filtering threshold retains 98--99\% of matches.}
\label{tab:s4}
\begin{tabular}{lccccc}
\toprule
City & $n$ & Median (m) & \% within 50\,m & \% within 100\,m & \% within 200\,m \\
\midrule
Jakarta  & 20,954 & 17.8 & 89.3\% & 97.7\% & 99.3\% \\
Bandung  & 4,382  & 19.6 & 86.3\% & 96.4\% & 98.9\% \\
Semarang & 1,539  & 19.6 & 85.2\% & 94.6\% & 97.9\% \\
\bottomrule
\end{tabular}
\end{table}

The high match quality ($>$85\% within 50\,m) confirms that the nearest-edge spatial join introduces minimal positional error, supporting the reliability of centrality--congestion correlation and capacity drop analyses that depend on this join.

%% ============================================================
\section{Spatial Regression Results}
\label{sec:s5}

To properly account for spatial dependence, Spatial Lag and Spatial Error models (via PySAL's \texttt{spreg}) were estimated alongside OLS multiple regression. The dependent variable is mean evening-peak jam factor; independent variables are edge betweenness centrality and POI density within a 500\,m buffer.

\begin{table}[!htbp]
\centering
\caption{Spatial regression results: jam factor $\sim$ betweenness centrality $+$ POI density. Across all model specifications, $R^2$ values are negligible ($<$0.3\%), confirming that spatial predictors explain virtually none of the congestion variance.}
\label{tab:s5}
\small
\begin{tabular}{llcccccl}
\toprule
City & Model & $R^2$ & $\beta$ (betw.) & $p$ & $\beta$ (POI) & $p$ & Spatial param. \\
\midrule
\multirow{3}{*}{Jakarta}
 & OLS           & 0.003\% & $-0.048$ & 0.635 & 0.000010  & 0.704 & ---              \\
 & Spatial Lag   & 0.006\% & $-0.337$ & 0.168 & $-0.000005$ & 0.899 & $\rho = 2.731$   \\
 & Spatial Error & 0.003\% & $-0.051$ & 0.615 & 0.000009  & 0.708 & $\lambda = 0.012$ \\
\midrule
\multirow{3}{*}{Bandung}
 & OLS           & 0.260\% & 0.270    & 0.005 & $-0.000010$ & 0.803 & ---              \\
 & Spatial Lag   & 0.128\% & 0.231    & 0.064 & $-0.000000$ & 0.996 & $\rho = 0.350$   \\
 & Spatial Error & 0.260\% & 0.269    & 0.005 & $-0.000010$ & 0.817 & $\lambda = 0.023$ \\
\midrule
\multirow{3}{*}{Semarang}
 & OLS           & 0.004\% & $-0.017$ & 0.894 & $-0.000007$ & 0.887 & ---              \\
 & Spatial Lag   & 0.013\% & 0.033    & 0.866 & 0.000055  & 0.706 & $\rho = 2.568$   \\
 & Spatial Error & 0.004\% & $-0.018$ & 0.886 & $-0.000008$ & 0.867 & $\lambda = -0.034$ \\
\bottomrule
\end{tabular}
\end{table}

Bandung's OLS betweenness coefficient is statistically significant ($p{=}0.005$) but trivially small in magnitude ($R^2 = 0.26\%$). This significance vanishes in the Spatial Lag model ($p{=}0.064$), suggesting it is partially attributable to spatial dependence rather than a genuine centrality effect. The spatial autoregressive parameters ($\rho$, $\lambda$) are small, consistent with the weak global spatial autocorrelation reported in the main text and Table~\ref{tab:s1}.

%% ============================================================
\section{Supplementary Figures}
\label{sec:sfigs}

\subsection{Hotspot Maps}

Figures~\ref{fig:s1}--\ref{fig:s3} show LISA cluster and Getis-Ord Gi* hotspot maps for the evening peak period across all three cities. The main text includes the Jakarta hotspot map (Figure~5); Bandung and Semarang maps are provided here.

\begin{figure}[!htbp]
\centering
\includegraphics[width=\textwidth]{figures/jkt_hotspots_evening_peak.png}
\caption{LISA cluster and Getis-Ord Gi* hotspot maps --- Jakarta (evening peak). Hot spots concentrate around the inner toll-road corridors and central business district, while cold spots cluster in peripheral residential areas.}
\label{fig:s1}
\end{figure}

\begin{figure}[!htbp]
\centering
\includegraphics[width=\textwidth]{figures/bdg_hotspots_evening_peak.png}
\caption{LISA cluster and Getis-Ord Gi* hotspot maps --- Bandung (evening peak). Hot spots are concentrated along the northern industrial corridor (Jl.\ Soekarno-Hatta) and the Pasteur--Cipaganti commercial strip.}
\label{fig:s2}
\end{figure}

\begin{figure}[!htbp]
\centering
\includegraphics[width=\textwidth]{figures/smg_hotspots_evening_peak.png}
\caption{LISA cluster and Getis-Ord Gi* hotspot maps --- Semarang (evening peak). Despite the smallest network ($n{=}1{,}076$), identifiable hot spots cluster around the Simpang Lima central roundabout and Jl.\ Pandanaran commercial corridor.}
\label{fig:s3}
\end{figure}

\subsection{Capacity--Congestion Scatter Plots}

\begin{figure}[!htbp]
\centering
\includegraphics[width=\textwidth]{figures/capacity_congestion_scatter.png}
\caption{Scatter plots of road capacity (lanes $\times$ speed limit) versus mean evening-peak jam factor for all three cities. Points show individual road segments; LOESS curves confirm the flat capacity--congestion relationship.}
\label{fig:s4}
\end{figure}

\subsection{Bottleneck Capacity Comparison}

\begin{figure}[!htbp]
\centering
\includegraphics[width=\textwidth]{figures/bottleneck_capacity_comparison.png}
\caption{Bottleneck capacity comparison across cities. Box plots compare the distribution of road capacity attributes (number of lanes, speed limits) between bottleneck and non-bottleneck segments identified by the graph-based capacity drop detection algorithm.}
\label{fig:s5}
\end{figure}

\subsection{Capacity Drop Spatial Distribution}

\begin{figure}[!htbp]
\centering
\includegraphics[width=\textwidth]{figures/capacity_drop_spatial_analysis.png}
\caption{Spatial distribution of capacity drop segments across all three cities. Segments are classified by whether they exhibit a statistically significant capacity drop (jam factor increase when transitioning from high-capacity to lower-capacity downstream segments).}
\label{fig:s6}
\end{figure}

\subsection{Street Orientation Analysis}

\begin{figure}[!htbp]
\centering
\includegraphics[width=\textwidth]{figures/street_orientation_polar.png}
\caption{Street orientation polar plots for Jakarta, Bandung, and Semarang. Jakarta shows a pronounced grid pattern with dominant N--S and E--W orientations; Bandung and Semarang show more heterogeneous orientations reflecting their topographic constraints.}
\label{fig:s7}
\end{figure}

%% ============================================================
\end{document}
